\documentclass[11pt]{article}
\usepackage[utf8]{inputenc}	% Para caracteres en español
\usepackage{amsmath,amsthm,amsfonts,amssymb,amscd}
\usepackage{multirow,booktabs}
\usepackage[table]{xcolor}
\usepackage{fullpage}
\usepackage{lastpage}
\usepackage{enumitem}
\usepackage{fancyhdr}
\usepackage{mathrsfs}
\usepackage{wrapfig}
\usepackage{setspace}
\usepackage{calc}
\usepackage{multicol}
\usepackage{cancel}
\usepackage{tikz}
\usepackage[retainorgcmds]{IEEEtrantools}
\usepackage[margin=3cm]{geometry}
\usepackage{amsmath}
\newlength{\tabcont}
\setlength{\parindent}{0.0in}
\setlength{\parskip}{0.05in}
\usepackage{empheq}
\usepackage{framed}
\usepackage{multicol}
\usepackage{amsmath}
\usepackage{setspace}
\usepackage{graphicx}
\usepackage{caption}
\usepackage{subcaption}
\usepackage{graphbox}
\graphicspath{ {figures/} }
\usepackage[most]{tcolorbox}
\usepackage{xcolor}
\colorlet{shadecolor}{orange!15}
\parindent 0in
\parskip 12pt
\geometry{margin=1in, headsep=0.25in}
\theoremstyle{definition}
\newtheorem{defn}{Definition}
\newtheorem{reg}{Rule}
\newtheorem{exer}{Exercise}
\newtheorem{note}{Note}
\newtheorem{example}{Example}
\newtheorem{theorem}{Theorem}
\newcommand{\indep}{\perp\!\!\!\!\perp} 
\begin{document}
\title{Report}

\thispagestyle{empty}

\begin{center}
{\LARGE \bf PHC 7091 Notes}\\
\end{center}
\tableofcontents
\newpage
\section{Generalized Linear Model}
\subsection{Canonical GLM}

Let $Y_1,\dots,Y_n$ be independent random variables. A canonical GLM assumes that
each $Y_i$ follows an exponential dispersion family distribution with density
\[
f(y_i \mid \theta_i, \phi)
=
\exp\!\left(
\frac{y_i \theta_i - b(\theta_i)}{a(\phi)}
+ c(y_i,\phi)
\right),
\]
where
\begin{itemize}
  \item $\theta_i$ is the \emph{natural (canonical) parameter},
  \item $\phi$ is the \emph{dispersion parameter},
  \item $a(\phi)$ is a known positive function,
  \item $b(\theta)$ is the cumulant (log-partition) function,
  \item $c(y,\phi)$ does not depend on $\theta_i$.
\end{itemize}
\paragraph{Canonical link.}
The canonical link function satisfies
\[
\theta_i = \eta_i = x_i^\top \beta,
\]
where $x_i$ is the covariate vector and $\beta$ is the regression coefficient vector.
\paragraph{Mean and variance.}
The first two moments are
\[
\mathbb{E}[Y_i] = \mu_i = b'(\theta_i),
\qquad
\mathrm{Var}(Y_i) = a(\phi)\, b''(\theta_i).
\]
Equivalently,
\[
\mathrm{Var}(Y_i) = a(\phi)\, V(\mu_i),
\quad
V(\mu_i) = b''(\theta(\mu_i)).
\]
\begin{center}
\renewcommand{\arraystretch}{1.4}
\begin{tabular}{lcccc}
\hline
Distribution
& $y$
& $\theta$
& $b(\theta)$
& $a(\phi)$
\\
\hline
Normal $(\mu,\sigma^2)$
& $y\in\mathbb{R}$
& $\mu$
& $\dfrac{\theta^2}{2}$
& $\sigma^2$
\\[6pt]

Poisson $(\lambda)$
& $y\in\{0,1,2,\dots\}$
& $\log\lambda$
& $e^{\theta}$
& $1$
\\[6pt]

Binomial $(n,p)$
& $y\in\{0,\dots,n\}$
& $\log\!\left(\dfrac{p}{1-p}\right)$
& $n\log(1+e^{\theta})$
& $1$
\\[6pt]

Bernoulli $(p)$
& $y\in\{0,1\}$
& $\log\!\left(\dfrac{p}{1-p}\right)$
& $\log(1+e^{\theta})$
& $1$
\\[6pt]

Gamma $(\mu,\phi)$
& $y>0$
& $-\dfrac{1}{\mu}$
& $-\log(-\theta)$
& $\phi$
\\[6pt]

Inverse Gaussian $(\mu,\phi)$
& $y>0$
& $-\dfrac{1}{2\mu^2}$
& $-\sqrt{-2\theta}$
& $\phi$
\\
\hline
\end{tabular}
\end{center}

\paragraph{Example}

For the Gamma GLM, Gamma is conventionally parameterized as
\[
Y_i \sim \mathrm{Gamma}\left(\frac{1}{\phi},\frac{1}{\phi\mu_i}\right),
\]
where
\[
\theta_i=-\frac{1}{\mu_i},\qquad
b(\theta_i)=-\log(-\theta_i),\qquad
a(\phi)=\phi,
\]
and
\[
c(y_i,\phi)
=
\left(\frac{1}{\phi}-1\right)\log y_i
-\frac{\log\phi}{\phi}
-\log\Gamma\left(\frac{1}{\phi}\right).
\]

\subsection{MLE $\theta$ for the Canonical Exponential Family}
Consider independent observations $Y_1,\dots,Y_n$ with density
\[
f(y_i \mid \theta_i,\phi)
=
\exp\!\left(
\frac{y_i \theta_i - b(\theta_i)}{a(\phi)}
+ c(y_i,\phi)
\right),
\]
where $a(\phi)>0$ is known.
\paragraph{Log-likelihood.}
The log-likelihood (up to constants) is
\[
\ell(\theta)
=
\frac{1}{a(\phi)}
\sum_{i=1}^n
\left(
y_i\theta_i - b(\theta_i)
\right).
\]
\paragraph{Score function.}
The score function with respect to $\theta_i$ is
\[
U(\theta_i)
=
\frac{\partial \ell}{\partial \theta_i}
=
\frac{1}{a(\phi)}
\left(
y_i - b'(\theta_i)
\right).
\]
\paragraph{Likelihood equation.}
The maximum likelihood estimator $\hat\theta_i$ satisfies
\[
b'(\hat\theta_i) = y_i,
\]
or equivalently,
\[
\hat\theta_i = (b')^{-1}(y_i),
\]
provided $b'(\theta)$ is invertible.
\paragraph{Existence and uniqueness.}
If $b(\theta)$ is strictly convex, then $b'(\theta)$ is strictly increasing,
and the MLE $\hat\theta_i$ exists and is unique whenever
$y_i$ lies in the interior of the mean parameter space.
\paragraph{Observed information.}
The observed Fisher information is
\[
-\frac{\partial^2 \ell}{\partial \theta_i^2}
=
\frac{1}{a(\phi)} b''(\theta_i).
\]

\paragraph{Fisher information.}
The Fisher information for $\theta_i$ is
\[
\mathcal{I}(\theta_i)
=
\mathbb{E}\!\left[
-\frac{\partial^2 \ell}{\partial \theta_i^2}
\right]
=
\frac{1}{a(\phi)} b''(\theta_i).
\]

\paragraph{Asymptotic distribution.}
Under standard regularity conditions,
\[
\sqrt{n}(\hat\theta - \theta)
\xrightarrow{d}
\mathcal{N}\!\left(
0,\,
\frac{a(\phi)}{b''(\theta)}
\right).
\]

\paragraph{Invariance property.}
Since $\mu = b'(\theta)$, the MLE of the mean parameter is
\[
\hat\mu = b'(\hat\theta).
\]

\paragraph{Efficiency.}
The MLE $\hat\theta$ attains the Cram\'er--Rao lower bound:
\[
\mathrm{Var}(\hat\theta) \ge \frac{a(\phi)}{n\,b''(\theta)},
\]
with equality asymptotically.

\paragraph{Newton update.}
Newton's method updates $\theta_i$ by
\[
\theta_i^{(t+1)}
=
\theta_i^{(t)}
-
\left[
\frac{\partial^2 \ell}{\partial \theta_i^2}
\right]^{-1}
\frac{\partial \ell}{\partial \theta_i}.
\]
Substituting the derivatives gives
\[
\boxed{
\theta_i^{(t+1)}
=
\theta_i^{(t)}
+
\frac{y_i - b'(\theta_i^{(t)})}{b''(\theta_i^{(t)})}
}
\]

\paragraph{Interpretation.}
The Newton step moves $\theta_i$ in proportion to
\[
\frac{\text{residual}}{\text{curvature}}
=
\frac{y_i - \mu_i^{(t)}}{\mathrm{Var}(Y_i)/a(\phi)},
\quad
\mu_i^{(t)} = b'(\theta_i^{(t)}).
\]

\paragraph{Convergence properties.}
\begin{itemize}
  \item If $b(\theta)$ is strictly convex, then $b''(\theta)>0$ and the update is well-defined.
  \item The log-likelihood is concave in $\theta$, so Newton iterations converge to the unique MLE.
  \item Convergence is quadratic in a neighborhood of $\hat\theta_i$.
\end{itemize}

\paragraph{Connection to Fisher scoring.}
Since
\[
\mathcal{I}(\theta_i)
=
-\mathbb{E}\!\left[
\frac{\partial^2 \ell}{\partial \theta_i^2}
\right]
=
\frac{1}{a(\phi)} b''(\theta_i),
\]
Newton's method coincides with Fisher scoring for canonical exponential families.

\paragraph{Special case (single observation).}
For a single observation $y$,
\[
\theta^{(t+1)}
=
\theta^{(t)}
+
\frac{y - b'(\theta^{(t)})}{b''(\theta^{(t)})},
\]
which converges to $\hat\theta = (b')^{-1}(y)$.

\subsection{IRLS with Mean $\mu_i$ and Link Function $g(\mu_i)$}

Let $Y_1,\dots,Y_n$ be independent with density
\[
f(y_i \mid \theta_i,\phi)
=
\exp\!\left(
\frac{y_i \theta_i - b(\theta_i)}{a(\phi)}
+ c(y_i,\phi)
\right),
\]
where
\[
\mu_i = \mathbb{E}[Y_i] = b'(\theta_i),
\qquad
\mathrm{Var}(Y_i) = a(\phi)\, b''(\theta_i).
\]

\paragraph{Link function.}
A generalized linear model assumes
\[
\eta_i = g(\mu_i) = x_i^\top \beta,
\]
where $g(\cdot)$ is a monotone link function.

\paragraph{Log-likelihood.}
Up to constants, the log-likelihood is
\[
\ell(\beta)
=
\frac{1}{a(\phi)}
\sum_{i=1}^n
\left(
y_i \theta_i - b(\theta_i)
\right),
\quad
\theta_i = \theta(\mu_i).
\]

\paragraph{Score function.}
Using the chain rule,
\[
\frac{\partial \ell}{\partial \beta}
=
\sum_{i=1}^n
\frac{\partial \ell}{\partial \mu_i}
\frac{\partial \mu_i}{\partial \eta_i}
\frac{\partial \eta_i}{\partial \beta}.
\]
Since
\[
\frac{\partial \ell}{\partial \mu_i}
=
\frac{y_i - \mu_i}{a(\phi)\, b''(\theta_i)},
\qquad
\frac{\partial \eta_i}{\partial \beta} = x_i,
\]
we obtain
\[
\nabla_\beta \ell(\beta)=
\frac{1}{a(\phi)}
\sum_{i=1}^n
\frac{y_i - b^{\prime}(\theta_i)}{b''(\theta_i)}
\frac{1}{g^\prime(\mu_i)}
x_i=
\frac{1}{a(\phi)}
\sum_{i=1}^n
\frac{y_i - \mu_i}{V(\mu_i)}
\frac{1}{g^\prime(\mu_i)}
x_i =
\frac{1}{a(\phi)}
\sum_{i=1}^n
\frac{y_i - \mu_i}{\text{Var}(\mu_i)}
\frac{1}{g^\prime(\mu_i)}
x_i
\]

\paragraph{Variance function.}
Define the variance function
\[
V(\mu_i) = b''(\theta_i),
\quad
\mathrm{Var}(Y_i) = a(\phi)\,V(\mu_i).
\]

\paragraph{Working weights.}
Define the weights
\[
w_i
=
\frac{1}{a(\phi)}
\left(\frac{d\mu_i}{d\eta_i}\right)^2
\frac{1}{V(\mu_i)} = \frac{1}{g^\prime(\mu_i)^2\text{Var}(Y_i)}\frac{1}{a(\phi)}
\]

\paragraph{Observed Hessian.}
Differentiating again gives the Hessian matrix
\[
\frac{\partial^2 l}{\partial \beta_k \partial \beta_j} = \frac{1}{a(\phi)}
\frac{\partial\sum_{i=1}^n
(y_i - \mu_i)
\frac{1}{g^\prime(\mu_i)V(\mu_i)}x_{ik}}{\partial\beta_j} 
\]
\[
= \frac{1}{a(\phi)}
\sum_{i=1}^n\frac{\partial (y_i - \mu_i)}{\partial\beta_j}
\frac{1}{g^\prime(\mu_i)V(\mu_i)} + (y_i - \mu_i)
\frac{\partial\frac{1}{g^\prime(\mu_i)V(\mu_i)}}{\partial\beta_j}x_{ik}
\]
\[
=\{\frac{1}{a(\phi)}\sum_{i=1}^n-\frac{\partial\mu_i}{\partial\eta_i}\frac{\partial\eta_i}{\partial\beta_i}\frac{1}{g^\prime(\mu_i)V(\mu_i)}-(y_i-\mu_i)(\frac{\partial\mu_i}{\partial\eta_i}\frac{\partial\eta_i}{\partial\beta_i}\frac{V^\prime(\mu_i)g^\prime(\mu_i) + V(\mu_i)g^{\prime\prime}(\mu_i)}{(V(\mu_i)g^{\prime}(\mu_i))^2})\}x_{ik}
\]
so with
\[
\frac{\partial\mu_i}{\partial\eta_j}=\frac{1}{g^\prime(\mu_i)}, \frac{\partial\eta_i}{\partial\beta_j}=x_{ij}
\]
we have it as 
\[
=\frac{1}{a(\phi)}\sum_{i=1}^n\{-\frac{1}{g^\prime(\mu_i)^2V(\mu_i)}-(y_i-\mu_i)(\frac{1}{g^\prime(\mu_i)}\frac{V^\prime(\mu_i)g^\prime(\mu_i) + V(\mu_i)g^{\prime\prime}(\mu_i)}{(V(\mu_i)g^{\prime}(\mu_i))^2})\}x_{ik}x_{ij}
\]
\paragraph{Matrix form.}
Define
\[
W_i^{\text{obs}}
=
\frac{1}{a(\phi)}
\left[
\frac{1}{V(\mu_i)}
\left(\frac{d\mu_i}{d\eta_i}\right)^2
-
(y_i-\mu_i)
\frac{d}{d\mu_i}
\left(
\frac{1}{V(\mu_i)}
\frac{d\mu_i}{d\eta_i}
\right)
\frac{d\mu_i}{d\eta_i}
\right],
\]
and let
\[
W^{\text{obs}} = \mathrm{diag}(W_1^{\text{obs}},\dots,W_n^{\text{obs}}).
\]
Then
\[
\boxed{
\nabla_\beta^2 \ell(\beta)
=
- X^\top W^{\text{obs}} X
}
\]
\subsection{IRLS for $\beta$ estimation}
\paragraph{Expected Hessian (Fisher information).}
Taking expectation with respect to $Y_i$ removes the residual term:
\[
\mathbb{E}[y_i-\mu_i]=0.
\]
Hence the Fisher information matrix is
\[
\boxed{
\mathcal{I}(\beta)
=
- \mathbb{E}[\nabla_\beta^2 \ell(\beta)]
=
X^\top W X
}
\]
with
\[
\boxed{
W_i
=
\frac{1}{a(\phi)}
\frac{1}{V(\mu_i)}
\left(\frac{d\mu_i}{d\eta_i}\right)^2
}
\]
\paragraph{Hessian (expected).}
The expected Hessian is
\[
\mathcal{I}(\beta)
=
X^\top W X,
\quad
W = \mathrm{diag}(w_1,\dots,w_n).
\]
\paragraph{Newton / Fisher scoring update.}
The Fisher scoring step is
\[
\beta^{(t+1)}
=
\beta^{(t)}
+
\left(X^\top W^{(t)} X\right)^{-1}
X^\top W^{(t)} (y - \mu^{(t)}),
\]
where all quantities are evaluated at $\beta^{(t)}$.

\paragraph{Working response.}
Define the working response
\[
z_i^{(t)}
=
\eta_i^{(t)}
+
(y_i - \mu_i^{(t)})
\left(\frac{d\eta_i}{d\mu_i}\right)_{\mu_i^{(t)}}.
\]

\paragraph{Weighted least squares form.}
Then $\beta^{(t+1)}$ solves
\[
\beta^{(t+1)}
=
\arg\min_{\beta}
\sum_{i=1}^n
w_i^{(t)}
\left(z_i^{(t)} - x_i^\top \beta\right)^2.
\]

\paragraph{IRLS algorithm.}
\begin{enumerate}
  \item Initialize $\beta^{(0)}$.
  \item At iteration $t$, compute
  \[
  \eta_i^{(t)} = x_i^\top \beta^{(t)}, \quad
  \mu_i^{(t)} = g^{-1}(\eta_i^{(t)}).
  \]
  \item Compute
  \[
  w_i^{(t)}
  =
  \frac{1}{a(\phi)}
  \left(\frac{d\mu_i}{d\eta_i}\right)^2
  \frac{1}{V(\mu_i)}.
  \]
  \item Form
  \[
  z_i^{(t)}
  =
  \eta_i^{(t)}
  +
  (y_i - \mu_i^{(t)})
  \left(\frac{d\eta_i}{d\mu_i}\right).
  \]
  \item Update
  \[
  \beta^{(t+1)}
  =
  \left(X^\top W^{(t)} X\right)^{-1}
  X^\top W^{(t)} z^{(t)}.
  \]
\end{enumerate}

\paragraph{Canonical link simplification.}
If $g(\mu)=\theta$, then
\[
\frac{d\mu}{d\eta} = V(\mu),
\quad
w_i = \frac{V(\mu_i)}{a(\phi)},
\]
and IRLS reduces to Newton's method exactly.
\subsection{Derivation of the IRLS Update via Taylor Expansion}

Assume a generalized linear model with
\[
\eta_i = g(\mu_i) = x_i^\top \beta,
\qquad
\mathrm{Var}(Y_i) = a(\phi)\,V(\mu_i),
\qquad
\mu_i = \mathbb{E}[Y_i].
\]

\paragraph{Score function.}
The score for $\beta$ is
\[
U(\beta)
=
\nabla_\beta \ell(\beta)
=
\sum_{i=1}^n
\frac{y_i - \mu_i}{a(\phi)\,V(\mu_i)}
\frac{d\mu_i}{d\eta_i}
x_i.
\]

\paragraph{First-order Taylor expansion of $\mu_i$.}
Expand $\mu_i(\beta)$ around the current iterate $\beta^{(t)}$:
\[
\mu_i(\beta)
\approx
\mu_i^{(t)}
+
\left.
\frac{d\mu_i}{d\eta_i}
\right|_{\beta^{(t)}}
x_i^\top (\beta - \beta^{(t)}).
\]

\paragraph{Linearization of the score.}
Substitute the expansion into the score equation $U(\beta)=0$:
\[
0
\approx
\sum_{i=1}^n
\frac{1}{a(\phi)\,V(\mu_i^{(t)})}
\frac{d\mu_i}{d\eta_i}
\left[
y_i
-
\mu_i^{(t)}
-
\frac{d\mu_i}{d\eta_i}
x_i^\top (\beta - \beta^{(t)})
\right]
x_i.
\]

\paragraph{Rearranging terms.}
Split the sum:
\[
\sum_{i=1}^n
\frac{y_i-\mu_i^{(t)}}{a(\phi)\,V(\mu_i^{(t)})}
\frac{d\mu_i}{d\eta_i}
x_i
=
\sum_{i=1}^n
\frac{1}{a(\phi)}
\frac{1}{V(\mu_i^{(t)})}
\left(\frac{d\mu_i}{d\eta_i}\right)^2
x_i x_i^\top
(\beta-\beta^{(t)}).
\]

\paragraph{Matrix form.}
Define the weights
\[
w_i^{(t)}
=
\frac{1}{a(\phi)}
\frac{1}{V(\mu_i^{(t)})}
\left(\frac{d\mu_i}{d\eta_i}\right)^2,
\qquad
W^{(t)} = \mathrm{diag}(w_1^{(t)},\dots,w_n^{(t)}).
\]
Then
\[
X^\top W^{(t)} X
(\beta-\beta^{(t)})
=
X^\top W^{(t)}
\left[
\eta^{(t)}
+
(y-\mu^{(t)})
\left(\frac{d\eta}{d\mu}\right)
-
X\beta^{(t)}
\right].
\]

\paragraph{Working response.}
Define the working response
\[
z_i^{(t)}
=
\eta_i^{(t)}
+
(y_i-\mu_i^{(t)})
\left(\frac{d\eta_i}{d\mu_i}\right).
\]

\paragraph{Normal equations.}
The update satisfies
\[
X^\top W^{(t)} X\,\beta^{(t+1)}
=
X^\top W^{(t)} z^{(t)}.
\]

\paragraph{IRLS update.}
Thus
\[
\boxed{
\beta^{(t+1)}
=
\left(X^\top W^{(t)} X\right)^{-1}
X^\top W^{(t)} z^{(t)}
}
\]

\paragraph{Canonical link simplification.}
If $g(\mu)=\theta$, then
\[
\frac{d\mu}{d\eta}=V(\mu),
\qquad
w_i^{(t)}=\frac{V(\mu_i^{(t)})}{a(\phi)},
\qquad
z_i^{(t)}=\eta_i^{(t)}+\frac{y_i-\mu_i^{(t)}}{V(\mu_i^{(t)})}.
\]
\subsection{Initial Values for IRLS}
Consider the model with only the intercept, then estimate $\beta_0$ with $\hat\beta_0$, set initial value as
$$
(0,0,\dots,\hat\beta_0)
$$
Example for Poisson GLM:

\textbf{Step 1: Likelihood}

Assume
\[
Y_i \sim \text{Poisson}(\mu_i), \qquad \log(\mu_i) = \beta_0.
\]
The Poisson density is
\[
P(Y_i=y_i) = \frac{\mu_i^{y_i} e^{-\mu_i}}{y_i!}.
\]
The log-likelihood for $n$ independent observations is
\[
l(\beta_0)
=
\sum_{i=1}^{n}
\left[
y_i \log(\mu_i) - \mu_i - \log(y_i!)
\right].
\]
Since $\mu_i = e^{\beta_0}$ and $\log(\mu_i)=\beta_0$,
\[
l(\beta_0)
=
\sum_{i=1}^{n}
\left[
y_i \beta_0 - e^{\beta_0} - \log(y_i!)
\right].
\]


\textbf{Step 2: Score Equation}

Differentiate the log-likelihood:
\[
\frac{\partial l(\beta_0)}{\partial \beta_0}
=
\sum_{i=1}^{n}
\left(
y_i - e^{\beta_0}
\right).
\]
Set the score equation equal to zero:
\[
\sum_{i=1}^{n} y_i - n e^{\beta_0} = 0.
\]


\textbf{Step 3: Maximum Likelihood Estimator}

Solving for $\beta_0$,
\[
e^{\hat{\beta}_0}
=
\frac{1}{n}\sum_{i=1}^{n} y_i
=
\bar{y}.
\]
Therefore
\[
\hat{\beta}_0
=
\log(\bar{y}).
\]
\textbf{Step 4: Update}
\[
\mu^{(0)} =
\bar y 
\quad \text{or} \quad
\frac{\bar y + y}{2}.
\]
Let
\[
\eta^{(0)} = X\beta^{(0)} + \text{offset}.
\]

Define the working response
\[
z
=
\eta^{(0)} - \text{offset}
+
\Delta\bigl(y-\mu^{(0)}\bigr),
\]
where
\[
\Delta = \frac{d\eta}{d\mu}.
\]
For the Poisson log-link,
\[
\eta = \log(\mu),
\qquad
\Delta = \frac{1}{\mu}.
\]
The weight matrix is
\[
W_0 = \operatorname{diag}(\mu^{(0)}).
\]
The first IRLS update is then
\[
\beta^{(1)}
\leftarrow
\left(X^\top W_0 X\right)^{-1}
X^\top W_0 z .
\]
\subsection{Equivalence of IRLS forms}
Using the working response
\[
z^{(t)}
=
\eta^{(t)}
+
\left.\frac{d\eta}{d\mu}\right|_{\mu^{(t)}}
\left(y-\mu^{(t)}\right),
\qquad
\eta^{(t)}=X\beta^{(t)},
\]
we have
\[
\begin{aligned}
\beta^{(t+1)}
&=
\left(X^\top W^{(t)}X\right)^{-1}
X^\top W^{(t)}z^{(t)}\\
&=
\left(X^\top W^{(t)}X\right)^{-1}
X^\top W^{(t)}
\left[
X\beta^{(t)}
+
D^{(t)}(y-\mu^{(t)})
\right]\\
&=
\beta^{(t)}
+
\left(X^\top W^{(t)}X\right)^{-1}
X^\top W^{(t)}D^{(t)}
(y-\mu^{(t)}),
\end{aligned}
\]
where
\[
D^{(t)}
=
\operatorname{diag}\left(
\left.\frac{d\eta_i}{d\mu_i}\right|_{\mu_i^{(t)}}
\right).
\]
Therefore,
\[
\beta^{(t+1)}
=
\beta^{(t)}
+
\left(X^\top W^{(t)}X\right)^{-1}
X^\top W^{(t)}(y-\mu^{(t)})
\]
holds only when \(D^{(t)}=I\), such as under the identity link, or when
\(D^{(t)}\) is absorbed into the definition of the weight matrix.
\newpage
\section{Residual Test}
\paragraph{Hat Matrix}
Consider the linear regression model
\[
Y = X\beta + \varepsilon, \qquad \varepsilon \sim (0,\sigma^2 I_n).
\]
The fitted values are
\[
\hat{Y} = X\hat{\beta},
\qquad 
\hat{\beta} = (X^\top X)^{-1}X^\top Y .
\]
Define the \emph{hat matrix}
\[
H = X (X^\top X)^{-1} X^\top .
\]
Then
\[
\hat{Y} = HY .
\]
The diagonal element $h_{ii}$ of $H$ is called the \emph{leverage} of observation $i$.
$$
\operatorname{Var}(e)=\sigma^2(I-H)
$$
\paragraph{Ordinary Residual}
The residual for observation $i$ is
\[
e_i = y_i - \hat{\mu}_i ,
\]
where
\[
\hat{\mu}_i = x_i^\top \hat{\beta}.
\]
\paragraph{Standardized Residual}
Let $\hat{\gamma}$ be an estimator of $\sigma$ (often $\hat{\sigma}$).  
The standardized residual is
\[
r_i =
\frac{y_i - \hat{\mu}_i}{\hat{\gamma}}.
\]
\[
\operatorname{Var}(e_i) = \sigma^2 (1 - h_{ii}).
\]
If $y_i\sim \operatorname{Poisson}(\mu_i)$. The standardized residual for observation $i$ is
\[
r_i
=
\frac{Y_i - \hat{\mu}_i}
{\sqrt{\hat{\mu}_i}} .
\]
\paragraph{Pearson Goodness-of-Fit Statistic}
Let $o_i$ denote the observed count and $e_i$ denote the expected count
under the fitted model, for $i=1,\dots,n$.

The Pearson goodness-of-fit statistic is
\[
X^2
=
\sum_{i=1}^n
\frac{(o_i - e_i)^2}{e_i}
\]
Under correct model specification and for large $n$,
\[
X^2 \approx \chi^2_{n-p},
\]
where $p$ is the number of estimated parameters in the model.
\[
\boxed{
X_P^2=\sum_j\frac{(O_j-E_j)^2}{E_j}
}
\]
This form is used when the observations are \textbf{cell counts or category
frequencies}, such as:
\begin{itemize}
    \item multinomial goodness-of-fit tests;
    \item contingency tables;
    \item Poisson count models.
\end{itemize}

For independent Poisson cell counts,
\[
O_j\sim\operatorname{Poisson}(E_j),
\]
so that
\[
E(O_j)=E_j,
\qquad
\operatorname{Var}(O_j)=E_j.
\]
Therefore,
\[
\frac{(O_j-E_j)^2}{E_j}
=
\frac{\left(O_j-E(O_j)\right)^2}
{\operatorname{Var}(O_j)},
\]
and hence
\[
X_P^2
=
\sum_j
\frac{\left(O_j-E(O_j)\right)^2}
{\operatorname{Var}(O_j)}.
\]
For multinomial or contingency-table data, the cell counts are not independent
because their total is fixed. However, a multinomial distribution can be obtained
by conditioning independent Poisson counts on their total. Therefore, the same
Pearson statistic is used.

For example, if
\[
Y\sim\operatorname{Binomial}(n,\pi),
\]
the two cell counts are
\[
O_1=Y,
\qquad
O_2=n-Y,
\]
with expected counts
\[
E_1=n\pi,
\qquad
E_2=n(1-\pi).
\]
Then
\[
\begin{aligned}
X_P^2
&=
\frac{(Y-n\pi)^2}{n\pi}
+
\frac{\left[(n-Y)-n(1-\pi)\right]^2}{n(1-\pi)}
\\
&=
\frac{(Y-n\pi)^2}{n\pi(1-\pi)}
\\
&=
\frac{\left(Y-E(Y)\right)^2}{\operatorname{Var}(Y)}.
\end{aligned}
\]
Thus,
\[
\boxed{
\begin{aligned}
\sum_i\frac{(Y_i-\mu_i)^2}{\operatorname{Var}(Y_i)}
&\quad\text{is used for general model responses,}\\
\sum_j\frac{(O_j-E_j)^2}{E_j}
&\quad\text{is used for cell-count or frequency data.}
\end{aligned}
}
\]
The denominator \(E_j\) is appropriate because cell counts are modeled directly
as Poisson counts, or as multinomial counts obtained by conditioning Poisson
counts on a fixed total.

\textbf{The connection between fixed total Poisson and Multinomial model will be provided in 2.3}
\subsection{Testing Nested Models in GLMs}
Let $M_0$ denote the model under the null hypothesis and $M_1$ the model under the alternative hypothesis, where $M_0$ is nested within $M_1$.

The procedure for testing nested models in a generalized linear model (GLM) is:
\begin{enumerate}
    \item Specify the null model $M_0$ and the alternative model $M_1$.
    \item Compute the goodness-of-fit statistics $G_0$ for $M_0$ and $G_1$ for $M_1$.
    \item Form a test statistic such as
    \[
    G_1 - G_0
    \quad \text{or} \quad
    \frac{G_1}{G_0}.
    \]
    \item Reject the null hypothesis based on the distribution of the chosen test statistic.
\end{enumerate}
In GLMs, the difference in deviances between nested models is commonly used and asymptotically follows a $\chi^2$ distribution.

The \textbf{Akaike Information Criterion (AIC)} is defined as
\[
AIC = -2\ell + 2p,
\]
where $\ell$ is the maximized log-likelihood and $p$ is the number of parameters.

The \textbf{Bayesian Information Criterion (BIC)} is defined as
\[
BIC = -2\ell + p\log(n),
\]
where $n$ is the sample size.

Models with smaller AIC or BIC values are preferred.

\paragraph{Scaled Deviance}
Let \(Y_1,\dots,Y_n\) be independent observations from an exponential
family distribution with log-likelihood
\[
\ell(\boldsymbol{\mu};\mathbf{y})
=
\sum_{i=1}^n
\frac{y_i\theta_i-b(\theta_i)}{\phi}
+
\sum_{i=1}^n c(y_i,\phi),
\]
where \(\theta_i\) is the canonical parameter corresponding to the mean
\(\mu_i\), and \(\phi\) is the dispersion parameter.

The scaled deviance is defined as
\[
D(\mathbf{y};\boldsymbol{\mu})
=
2\left[
\ell(\widehat{\boldsymbol{\mu}}^{\mathrm{sat}};\mathbf{y})
-
\ell(\boldsymbol{\mu};\mathbf{y})
\right],
\]
where
\begin{itemize}
    \item \(\ell(\boldsymbol{\mu};\mathbf{y})\) is the log-likelihood
    evaluated at the fitted mean vector
    \(\boldsymbol{\mu}=(\mu_1,\dots,\mu_n)\);
    \item \(\widehat{\boldsymbol{\mu}}^{\mathrm{sat}}\) is the fitted mean
    vector under the saturated model, which \textbf{assigns a separate mean
    parameter to every observation}. When allowed by the mean parameter
    space, this gives
    \[
    \widehat{\mu}_i^{\mathrm{sat}}=y_i.
    \]
\end{itemize}

Thus, the deviance measures twice the difference between the log-likelihood
of the saturated model and that of the fitted model. It quantifies the loss
of fit resulting from imposing the structure of the fitted GLM rather than
allowing each observation to have its own mean parameter.

\subsection{Residual test examples}

\textbf{Pearson residual (Poisson case):}
\[
r_i =
\frac{y_i - \hat{\mu}_i}{\sqrt{\hat{\mu}_i}}
\]
\textbf{Pearson residual (Gamma case):}
\[
r_i =
\frac{y_i - \hat{\mu}_i}{\hat{\mu}_i}
\]
Under the model assumptions,
\[
\sum_{i=1}^n r_i^2 \sim \chi^2_{n-p}.
\]
\textbf{Deviance:}
\[
D(y,\hat{\mu}) =
\sum_{i=1}^n
2\phi
\left[
\ell(y_i;y_i) - \ell(\hat{\mu}_i;y_i)
\right].
\]

\textbf{Deviance residual:}
\[
r_i^d =
\operatorname{sign}(y_i-\hat{\mu}_i)
\sqrt{d_i},
\]
where $d_i$ is the individual contribution to the deviance.

The scaled deviance satisfies
\[
\frac{D(y,\hat{\mu})}{\phi} \sim \chi^2_{n-p}.
\]
\textbf{Null Deviance.}

The \emph{null deviance} is obtained by fitting the \textbf{null model}, which contains only an intercept (no covariates).  
It is computed by comparing the saturated model with this intercept-only model.
\paragraph{Likelihood Ratio Test for Nested GLMs}
Let $M_0 \subset M_1$ be two nested GLMs with $p_0 < p_1$ parameters.  
Let $D(y,\hat{\mu}_0)$ and $D(y,\hat{\mu}_1)$ denote the deviances under $M_0$ and $M_1$.

The likelihood ratio test statistic is
\[
G
=
\frac{D(y,\hat{\mu}_0)-D(y,\hat{\mu}_1)}{\phi} = 2(\ell(y,\hat\mu_{full}) - \ell(y,\hat\mu_{reduced}))\sim \chi^2_{p.full-p.reduced}
\]
$\phi$ can be estimated from the full model

\paragraph{Example: Poisson Regression}
Assume the number of deaths $Y$ follows a Poisson distribution
\[
f(y;\mu)=\frac{\mu^y e^{-\mu}}{y!}.
\]

The expected number of deaths depends on the population size $M$ and a mortality rate that depends on covariates $x$. Thus we model
\[
E(Y) = \mu = M \mathcal{X}(x^\prime \beta),
\]
where $\mathcal{X}(x^\prime \beta)$ represents the mortality rate.

A common specification (Dobson) is
\[
E(Y) = \mu = M e^{x^\prime \beta}.
\]

Using the log link,
\[
g(\mu_i) = \log(\mu_i) = \log(M_i) + x_i^\prime \beta,
\]
where $\log(M_i)$ acts as an offset.

Detail at textbook 9.2, P198
$$
E\left(Y_i\right)=\mu_i=n_i \theta_i
$$
$$
\theta_i=e^{\mathbf{x}_i^T \boldsymbol{\beta}}
$$
The GLM model is
$$
E\left(Y_i\right)=\mu_i=n_i e^{\mathbf{x}_i^T \boldsymbol{\beta}} ; \quad Y_i \sim \operatorname{Po}\left(\mu_i\right)
$$
Link:
$$
\log \mu_i=\log n_i (\text{offset})+\mathbf{x}_i^T \boldsymbol{\beta}
$$
and we have rate ratio
$$
R R=\frac{\mathrm{E}\left(Y_i \mid \text { present }\right)}{\mathrm{E}\left(Y_i \mid \text { absent }\right)}=e^{\beta_j}
$$
All the other explanatory variables remain the same, for a continuous explanatory variable $x_k$, a one-unit increase will result in a multiplicative effect of $e^{\beta_k}$ on the rate $\mu$.

For a binary explanatory variable denoted by an indictor variable, $x_j=0$ if the factor is absent and $x_j=1$ if it is present. 
$$
\frac{b_j-\beta_j}{\text { s.e. }\left(b_j\right)} \sim N(0,1)
$$
where $b_j=\hat\beta_j$ is the estimator. For Wald test, confidence intervals, score, LRT, etc.
Fit 
$$
\widehat{Y}_i=\widehat{\mu}_i=n_i e^{\mathbf{x}_i^T \mathbf{b}}
$$
Also denoted as $e_i$ ($e_i = \hat Y_i$)

Pearson residuals
$$
r_i=\frac{o_i-e_i}{\sqrt{e_i}} = \frac{Y_i-\hat Y_i}{\sqrt{\hat Y_i}}
$$
$o_i$ denotes the observed value of $Y_i$. The deviance is given by
$$
D=2 \sum\left[o_i \log \left(o_i / e_i\right)-\left(o_i-e_i\right)\right]
$$
However, for most models $\sum o_i=\sum e_i$, so the deviance simplifies to
$$
D=2 \sum\left[o_i \log \left(o_i / e_i\right)\right] .
$$
With Tarlor expansion
$$
o \log \left(\frac{o}{e}\right)=(o-e)+\frac{1}{2} \frac{(o-e)^2}{e}+\dots
$$
we have
$$
\begin{aligned}
D & =2 \sum_{i=1}^N\left[\left(o_i-e_i\right)+\frac{1}{2} \frac{\left(o_i-e_i\right)^2}{e_i}-\left(o_i-e_i\right)\right] \\
& =\sum_{i=1}^N \frac{\left(o_i-e_i\right)^2}{e_i}=X^2=\sum r_i^2.
\end{aligned}
$$
Both $D$ and $X^2$ follows $\chi^2_{N-p}$. where \(N\) is the number of independent observations and \(p\) is the
number of estimated regression parameters, including the intercept. The
offset is known and is therefore not counted as an estimated parameter.

Example 9.2.1, P201
\subsection{Connection between fixed total Poisson and Multinomial model}

\paragraph{Example: Multinomial Model}
If the only constraint is that the sum of the $Y_i$ 's is $n$, then the following Multinomial distribution may be used
$$
f(\mathbf{y} ; \boldsymbol{\mu} \mid n)=n!\prod_{i=1}^N \theta_i^{y_i} / y_{i}!,
$$
where $\sum_{i=1}^N \theta_i=1$ and $\sum_{i=1}^N y_i=n$. In this case, $E\left(Y_i\right)=n \theta_i$.

\paragraph{Example: Poisson Model with Different rate}
Let $Y_1,\dots,Y_T$ be independent random variables such that
\[
Y_i \sim \text{Poisson}(\lambda_i), \qquad i=1,\dots,T .
\]
The probability mass function of $Y_i$ is
\[
P(Y_i=y_i)
=
\frac{e^{-\lambda_i}\lambda_i^{y_i}}{y_i!}.
\]
Given observations $y_1,\dots,y_T$, the likelihood function is
\[
L(\lambda_1,\dots,\lambda_T; y_1,\dots,y_T)
=
\prod_{i=1}^{T}
\frac{e^{-\lambda_i}\lambda_i^{y_i}}{y_i!}.
\]
that
$$
\sum Y_i \sim \operatorname{Poisson} (\sum \lambda_i)
$$
The log-likelihood is
\[
\ell(\lambda_1,\dots,\lambda_T)
=
\sum_{i=1}^{T}
\left(
-\lambda_i + y_i\log\lambda_i - \log(y_i!)
\right).
\]
\paragraph{Poisson and Multinomial Conversion}
Let $Y_1,\dots,Y_T$ be independent random variables with
\[
Y_i \sim \operatorname{Poisson}(\lambda_i), \qquad i=1,\dots,T .
\]
Then
\[
S=\sum_{i=1}^T Y_i \sim \operatorname{Poisson}\!\left(\sum_{i=1}^T \lambda_i\right).
\]
Moreover, conditional on $S=s$,
\[
(Y_1,\dots,Y_T)\mid S=s
\sim
\operatorname{Multinomial}
\left(s;\,
\frac{\lambda_1}{\sum_{i=1}^T\lambda_i},
\dots,
\frac{\lambda_T}{\sum_{i=1}^T\lambda_i}
\right).
\]

\begin{proof}
Let $S=\sum_{i=1}^T Y_i$. Since the $Y_i$ are independent,
\[
P(Y_1=y_1,\dots,Y_T=y_T)
=
\prod_{i=1}^T
\frac{e^{-\lambda_i}\lambda_i^{y_i}}{y_i!}.
\]
To compute the distribution of $S$, write
\[
P(S=s)
=
\sum_{\substack{y_1+\cdots+y_T=s \\ y_i\ge0}}
P(Y_1=y_1,\dots,Y_T=y_T).
\]
Substituting the joint pmf gives
\[
P(S=s)
=
\sum_{y_1+\cdots+y_T=s}
\prod_{i=1}^T
\frac{e^{-\lambda_i}\lambda_i^{y_i}}{y_i!}.
\]
Factor out the terms not depending on $(y_1,\dots,y_T)$:
\[
P(S=s)
=
e^{-\sum_{i=1}^T\lambda_i}
\sum_{y_1+\cdots+y_T=s}
\prod_{i=1}^T
\frac{\lambda_i^{y_i}}{y_i!}.
\]
Use the multinomial expansion
\[
(\lambda_1+\cdots+\lambda_T)^s
=
\sum_{y_1+\cdots+y_T=s}
\frac{s!}{y_1!\cdots y_T!}
\prod_{i=1}^T \lambda_i^{y_i}.
\]
Hence
\[
\sum_{y_1+\cdots+y_T=s}
\prod_{i=1}^T
\frac{\lambda_i^{y_i}}{y_i!}
=
\frac{(\sum_{i=1}^T\lambda_i)^s}{s!}.
\]
Therefore
\[
P(S=s)
=
e^{-\sum_{i=1}^T\lambda_i}
\frac{(\sum_{i=1}^T\lambda_i)^s}{s!},
\]
which is the pmf of
\[
\operatorname{Poisson}\!\left(\sum_{i=1}^T\lambda_i\right).
\]
Now compute the conditional distribution. For $y_1+\cdots+y_T=s$,
\[
P(Y_1=y_1,\dots,Y_T=y_T\mid S=s)
=
\frac{P(Y_1=y_1,\dots,Y_T=y_T)}{P(S=s)}.
\]
Substituting the expressions above and simplifying gives
\[
P(Y_1=y_1,\dots,Y_T=y_T\mid S=s)
=
\frac{s!}{y_1!\cdots y_T!}
\prod_{i=1}^T
\left(
\frac{\lambda_i}{\sum_{j=1}^T\lambda_j}
\right)^{y_i}.
\]
This is the pmf of a multinomial distribution with parameters
\[
s
\quad\text{and}\quad
p_i=\frac{\lambda_i}{\sum_{j=1}^T\lambda_j}.
\]
\end{proof}
Therefore, conditioning on $S=s$, since
\[
(Y_1,\dots,Y_T)\mid S=s
\sim
\operatorname{Multinomial}(s; p_1,\dots,p_T),
\qquad
p_i=\frac{\lambda_i}{\sum_{j=1}^T\lambda_j},
\]
we have
\[
E(Y_i\mid S=s)=sp_i,
\]
\[
\operatorname{Var}(Y_i\mid S=s)=sp_i(1-p_i),
\]
and for $i\neq j$,
\[
\operatorname{Cov}(Y_i,Y_j\mid S=s)=-sp_ip_j.
\]
Log-linear model is often used for associations.

Example 9.7.2 P215, 9.3.2 on P208

7.2, 7.3, 7.4
\newpage

\section{Binomial Model GLM}
\paragraph{Binomial GLM}
Let $Y_1,\dots,Y_N$ be independent random variables such that
\[
Y_i \sim \operatorname{Binomial}(\mu_i,\pi_i), \qquad i=1,\dots,N,
\]
where $\mu_i$ is the number of trials and $\pi_i$ is the success probability.

The generalized linear model specifies that the mean
\[
E(Y_i)=\mu_i\pi_i
\]
is linked to the covariates through
\[
g(\pi_i)=x_i^{\prime}\beta,
\]
where $g(\cdot)$ is the link function.

The likelihood function is
\[
L(\pi;y)
=
\prod_{i=1}^{N}
\binom{\mu_i}{y_i}
\pi_i^{y_i}(1-\pi_i)^{\mu_i-y_i}.
\]
The log-likelihood is
\[
\ell(\pi;y)
=
\sum_{i=1}^{N}
\left[
y_i\log(\pi_i)
+
(\mu_i-y_i)\log(1-\pi_i)
+
\log\binom{\mu_i}{y_i}
\right].
\]
The maximum likelihood estimator $\hat{\beta}$ maximizes $\ell(\pi;y)$ subject to
\[
g(\pi_i)=x_i^{\prime}\beta.
\]
Table 7.1
\paragraph{Deviance for the Binomial GLM}
Let $Y_1,\dots,Y_N$ be independent with
\[
Y_i \sim \operatorname{Binomial}(\mu_i,\pi_i),
\qquad
E(Y_i)=\mu_i\pi_i .
\]
Let $\hat{\pi}_i$ denote the fitted probabilities from the GLM and
$\hat{\mu}_i=\mu_i\hat{\pi}_i$ the fitted means.

The deviance is defined as
\[
D(y;\hat{\mu})
=
2\bigl[l(y;y)-l(\hat{\mu};y)\bigr].
\]

For the binomial model this becomes
\[
D(y;\hat{\mu})
=
2\sum_{i=1}^{N}
\left[
y_i\log\!\left(\frac{y_i}{\hat{\mu}_i}\right)
+
(\mu_i-y_i)\log\!\left(
\frac{\mu_i-y_i}{\mu_i-\hat{\mu}_i}
\right)
\right].
\]
That is
\[
D(y;\hat{\mu})
=
2\sum_{i=1}^{N}
\left[
y_i\log\!\left(\frac{y_i}{\hat{\mu}_i}\right)
\right].
\]
that
\[
D = 2\sum o\log(o/e)
\]
that $o$ is the observation ($y_i$ or $y_i -\mu_i$), $e$ is the respective frequency (fitted values). $\widehat{y}_i=n_i \widehat{\pi}_i$ and $\left(n_i-\widehat{y}_i\right)=\left(n_i-n_i \widehat{\pi}_i\right)$
that
$$
D\sim\chi^2_{N-p}
$$
Suppose model 1 with $p_1$ parameters and model 2 with $p_2$ parameters $p_1 < p_2$ that
\[
\Delta D = D_1 - D_2 \sim \chi^2_{p_2-p_1}
\]
This could be taken for examine if we prefer model 1 or 2. We can also use AIC

(7.4 P158)

(7.5, P162)
\paragraph{Binomial GLM Goodness of fit}
$$
S_w=\sum_{i=1}^N \frac{\left(y_i-n_i \pi_i\right)^2}{n_i \pi_i\left(1-\pi_i\right)}
$$
since $E\left(Y_i\right)=n_i \pi_i$ and $\operatorname{Var}\left(Y_i\right)=n_i \pi_i\left(1-\pi_i\right)$.
This is equivalent to minimizing the Pearson chi-squared statistic
$$
X^2=\sum \frac{(o-e)^2}{e},
$$
\begin{proof}
    $$
\begin{aligned}
X^2 & =\sum_{i=1}^N \frac{\left(y_i-n_i \pi_i\right)^2}{n_i \pi_i}+\sum_{i=1}^N \frac{\left[\left(n_i-y_i\right)-n_i\left(1-\pi_i\right)\right]^2}{n_i\left(1-\pi_i\right)} \\
& =\sum_{i=1}^N \frac{\left(y_i-n_i \pi_i\right)^2}{n_i \pi_i\left(1-\pi_i\right)}\left(1-\pi_i+\pi_i\right)=S_w .
\end{aligned}
$$
\end{proof}
$X^2$ is estimated at
$$
X^2=\sum_{i=1}^N \frac{\left(y_i-n_i \widehat{\pi}_i\right)^2}{n_i \widehat{\pi}_i\left(1-\widehat{\pi}_i\right)}
$$
which is asymptotically equivalent to the deviance
with taylor expansion
$$
o\log(o/e) = (o-e) + \frac{1}{2}\frac{(o-e)^2}{e}+\dots
$$
that gives
$$
\begin{aligned}
D = & 2 \sum_{i=1}^N\left\{\left(y_i-n_i \widehat{\pi}_i\right)+\frac{1}{2} \frac{\left(y_i-n_i \widehat{\pi}_i\right)^2}{n_i \widehat{\pi}_i}+\left[\left(n_i-y_i\right)-\left(n_i-n_i \widehat{\pi}_i\right)\right]\right. \\
& \left.+\frac{1}{2} \frac{\left[\left(n_i-y_i\right)-\left(n_i-n_i \widehat{\pi}_i\right)\right]^2}{n_i-n_i \widehat{\pi}_i}+\ldots\right\} \\
\cong & \sum_{i=1}^N \frac{\left(y_i-n_i \widehat{\pi}_i\right)^2}{n_i \widehat{\pi}_i\left(1-\widehat{\pi}_i\right)}=X^2 
\end{aligned}
$$
Another goodness of fit testing is 
$$
C=2[l(\widehat{\boldsymbol{\pi}} ; \mathbf{y})-l(\widetilde{\boldsymbol{\pi}} ; \mathbf{y})]
$$
for test statistics, where $l$ is the log-likelihood function
$$
C=2 \sum\left[y_i \log \left(\frac{\widehat{y}_i}{n_i \widetilde{\pi}_i}\right)+\left(n_i-y_i\right) \log \left(\frac{n_i-\widehat{y}_i}{n_i-n_i \widetilde{\pi}_i}\right)\right]
$$
Under the minimal model $\widetilde{\pi}=\left(\Sigma y_i\right) /\left(\Sigma n_i\right)$. Let $\widehat{\pi}_i$ denote the estimated probability for $Y_i$ under the model of interest (so the fitted value is $\left.\widehat{y}_i=n_i \widehat{\pi}_i\right)$.

This tests if all $\beta$ entries are 0 or at least one of them is non-zero. the approximate sampling distribution for $C$ is $\chi^2_{p-1}$ if all the $p$ parameters except the intercept term are zero

(7.6 P166)

The Pearson, or chi-squared, residual is
$$
X_k=\frac{\left(y_k-n_k \widehat{\pi}_k\right)}{\sqrt{n_k \widehat{\pi}_k\left(1-\widehat{\pi}_k\right)}} \quad, k=1, \ldots, m .
$$
Deviance residuals can be defined similarly,
$$
d_k=\operatorname{sign}\left(y_k-n_k \widehat{\pi}_k\right)\left\{2\left[y_k \log \left(\frac{y_k}{n_k \widehat{\pi}_k}\right)+\left(n_k-y_k\right) \log \left(\frac{n_k-y_k}{n_k-n_k \widehat{\pi}_k}\right)\right]\right\}^{1 / 2}
$$
Both provides diagnotic information
\begin{itemize}
    \item Residuals vs each explanatory variable (check for linearity)
    \item Plotted in the order of the measurements (check for serial correlation)
    \item Normal probability plots (Assume standardized residual should have approximately normal)
\end{itemize}
\paragraph{Example: Does response model}
We have the $y_i$ as the success out of $n_i$ trials, with the proportion of success as $p_i = y_i/n_i$. We have $E(p_i) = \pi_i$ and model $g(\pi_i)=x_i^\prime\beta$ as the link function
\[
f(y;\pi)
=
\binom{n}{np}
\pi^{np}(1-\pi)^{n-np} = 
\exp\!\left\{
\frac{p\log\frac{\pi}{1-\pi}
+
\log(1-\pi)}{1/n}+ \log\binom{n}{np}
\right\}.
\]
with $\theta = \log(\pi/(1-\pi))$, it becomes
\[
f(p;\pi)
=
\binom{n}{np}
\exp\!\left\{
np\,\theta
-
n\log(1+e^{\theta})
\right\},
\]
\[
\log f(p;\pi)
=
\log\binom{n}{np}
+
\frac{p\,\theta
-
\log(1+e^{\theta})}{1/n}
\]
We have
$$
\pi = \frac{e^\theta}{1+e^\theta} = E(p)
$$
and
$$
\operatorname{Var}(p_i)
=
\frac{a(\phi)}{n}
\frac{e^{\theta_i}}{(1+e^{\theta_i})^2}.
$$
with the link function
\[
g(\pi_i) = \eta_i = \beta_1 + \beta_2x_i
\]
\subsection{IRLS for Binomial GLM}
Let $Y_i \sim \text{Binomial}(n_i,\pi_i)$ and define the observed
proportions
\[
p_i=\frac{y_i}{n_i}, \qquad
p=(p_1,\dots,p_n)^\prime .
\]

Assume the GLM
\[
g(\pi_i)=x_i^\prime \beta, \qquad
\eta = X\beta .
\]

The MLE of $\beta$ can be obtained using the
\textbf{Iteratively Reweighted Least Squares (IRLS)} algorithm.
At iteration $m$ we solve the weighted least squares equation
\[
(X^\prime W X)^{(m-1)} b^{(m)} = X^\prime W z .
\]
The components are defined as follows.

\textbf{Weight matrix}
\[
W
=
\operatorname{diag}\!\left(
\frac{1}{\operatorname{Var}(p_i)\,[g'(\pi_i)]^2}
\right)
=
\operatorname{diag}\!\left(
\frac{n_i}{\pi_i(1-\pi_i)\,[g'(\pi_i)]^2}
\right).
\]

\textbf{Adjusted response}
\[
z
=
X b^{(m-1)}
+
\Delta (p-\pi)
=
\eta + \Delta (p-\pi),
\]
evaluated at $\beta=b^{(m-1)}$.

\textbf{Derivative matrix}
\[
\Delta = \operatorname{diag}\!\left(g'(\pi_i)\right).
\]

\textbf{Initial estimate}
\[
\hat b^{(0)} = (X^\prime X)^{-1}X^\prime y .
\]

The procedure iterates until convergence of $b^{(m)}$.

\subsection*{Example: IRLS for Binomial GLM}
Let
\[
Y_i \sim \text{Binomial}(n_i,\pi_i), 
\qquad 
p_i=\frac{y_i}{n_i},
\qquad 
p=(p_1,\dots,p_n)' .
\]

Assume
\[
\eta_i = \beta_1 + \beta_2 x_i,
\qquad 
\pi_i = \phi(\eta_i).
\]

The MLE of $\beta=(\beta_1,\beta_2)'$ can be obtained by the
\textbf{Iteratively Reweighted Least Squares (IRLS)} algorithm.
At iteration $m$ we solve
\[
(X^\prime W X)^{(m-1)} b^{(m)} = X^\prime W z .
\]
\textbf{Weight matrix}
\[
W
=
\operatorname{diag}\!\left(
\frac{1}{\operatorname{Var}(p_i)\,[\phi'(\eta_i)]^2}
\right)
=
\operatorname{diag}\!\left(
\frac{n_i[\phi'(\eta_i)]^2}{\pi_i(1-\pi_i)}
\right).
\]
\textbf{Adjusted response}
\[
z
=
X b^{(m-1)}
+
\Delta (p-\pi)
=
\eta + \Delta (p-\pi),
\]
where
\[
\Delta
=
\operatorname{diag}\!\left(
\frac{1}{\phi'(\eta_i)}
\right).
\]
\textbf{Design matrix}
\[
X=
\begin{pmatrix}
1 & x_1\\
1 & x_2\\
\vdots & \vdots\\
1 & x_n
\end{pmatrix}.
\]
The iteration proceeds until $b^{(m)}$ converges.

\subsection*{Example: Logistic Regression Model}
Let
\[
Y_i \sim \operatorname{Binomial}(n_i,\pi_i), 
\qquad 
p_i=\frac{y_i}{n_i}.
\]

The logistic link function is
\[
g(\pi_i)=\log\!\left(\frac{\pi_i}{1-\pi_i}\right)=\eta_i,
\]
with linear predictor
\[
\eta_i=\beta_1+\beta_2 x_i .
\]
Equivalently,
\[
\pi_i=\frac{e^{\eta_i}}{1+e^{\eta_i}}.
\]
The maximum likelihood estimate of $\beta=(\beta_1,\beta_2)'$
can be obtained using the \textbf{Iteratively Reweighted Least Squares (IRLS)} algorithm,
which at iteration $m$ solves
\[
(X'WX)^{(m-1)} b^{(m)} = X'Wz .
\]
The matrices are
\[
W=\operatorname{diag}\!\left(n_i\pi_i(1-\pi_i)\right),
\]
\[
\Delta=\operatorname{diag}(g^\prime(\pi_i)) = \operatorname{diag}(\frac{1}{\pi_i(1-\pi_i)})
\]
\[
z=\eta+\frac{p-\pi}{\pi(1-\pi)},
\]
and
\[
X=
\begin{pmatrix}
1 & x_1\\
1 & x_2\\
\vdots & \vdots\\
1 & x_n
\end{pmatrix}.
\]
The iteration proceeds until convergence of $b^{(m)}$.

\subsection*{Example: Binomial GLM with inverse log–log link}
Let
\[
Y_i \sim \operatorname{Binomial}(n_i,\pi_i), 
\qquad 
p_i=\frac{y_i}{n_i}.
\]
The link function is
\[
g(\pi_i)=\log\!\big(-\log(1-\pi_i)\big)=\eta_i,
\qquad
\eta_i=\beta_1+\beta_2 x_i .
\]
Hence
\[
\pi_i
=
1-\exp\!\big(-e^{\eta_i}\big).
\]
The MLE of $\beta=(\beta_1,\beta_2)'$ can be obtained by
\textbf{Iteratively Reweighted Least Squares (IRLS)}
\[
(X'WX)^{(m-1)} b^{(m)} = X'Wz .
\]
\textbf{Derivative of the inverse link}
\[
\phi'(\eta_i)
=
\frac{d\pi_i}{d\eta_i}
=
e^{\eta_i}\exp(-e^{\eta_i}).
\]
\textbf{Derivative of the link}
\[
g'(\pi_i)
=
\frac{1}{(1-\pi_i)\,[-\log(1-\pi_i)]}.
\]
\textbf{Weight matrix}
\[
W
=
\operatorname{diag}\!\left(
\frac{1}{\operatorname{Var}(p_i)\,[g'(\pi_i)]^2}
\right)
=
\operatorname{diag}\!\left(
n_i
\frac{(1-\pi_i)^2[\log(1-\pi_i)]^2}{\pi_i(1-\pi_i)}
\right).
\]
\textbf{Derivative matrix}
\[
\Delta
=
\operatorname{diag}\!\big(g'(\pi_i)\big)
=
\operatorname{diag}\!\left(
\frac{1}{(1-\pi_i)\,[-\log(1-\pi_i)]}
\right).
\]
\textbf{Adjusted response}
\[
z
=
X b^{(m-1)} + \Delta(p-\pi)
=
\eta + \Delta(p-\pi).
\]

\subsection*{Example: Binomial GLM with log–log link}
Let
\[
Y_i \sim \operatorname{Binomial}(n_i,\pi_i),
\qquad
p_i=\frac{y_i}{n_i}.
\]
Assume the link
\[
g(\pi_i)=\log(-\log \pi_i)=\eta_i,
\qquad
\eta_i=\beta_1+\beta_2 x_i .
\]
Hence the inverse link is
\[
\pi_i
=
\exp(-e^{\eta_i}).
\]
The MLE of $\beta$ can be obtained by the
\textbf{Iteratively Reweighted Least Squares (IRLS)} algorithm
\[
(X'WX)^{(m-1)} b^{(m)} = X'Wz .
\]
\textbf{Variance of the proportion}
\[
\operatorname{Var}(p_i)=\frac{\pi_i(1-\pi_i)}{n_i}.
\]
\textbf{Derivative of the link}
\[
g'(\pi_i)
=
-\frac{1}{\pi_i \log \pi_i}.
\]
\textbf{Weight matrix}
\[
W
=
\operatorname{diag}\!\left(
\frac{1}{\operatorname{Var}(p_i)\,[g'(\pi_i)]^2}
\right)
=
\operatorname{diag}\!\left(
\frac{n_i \pi_i (\log \pi_i)^2}{1-\pi_i}
\right).
\]
\textbf{Derivative matrix}
\[
\Delta
=
\operatorname{diag}\!\big(g'(\pi_i)\big)
=
\operatorname{diag}\!\left(
-\frac{1}{\pi_i \log \pi_i}
\right).
\]
\textbf{Adjusted response}
\[
z
=
X b^{(m-1)}+\Delta(p-\pi)
=
\eta+\Delta(p-\pi).
\]
Example 7.3
\newpage
\section{Multinomial logistic model}
\subsection{Logistic Model with Binary Covariate}
Suppose the response variable $Y$ has $J$ categories with probabilities
\[
\pi_j = P(Y=j), \qquad j=1,\dots,J,
\]
and let $x \in \{0,1\}$ be a binary covariate.

Taking category $1$ as the baseline, the multinomial logit model is
\[
\log\left(\frac{\pi_j}{\pi_1}\right)
=
\beta_{0j} + \beta_{1j}x,
\qquad j=2,\dots,J.
\]

\paragraph{Interpretation for $x=0$ and $x=1$}
When $x=0$,
\[
\log\left(\frac{\pi_{j0}}{\pi_{10}}\right)
=
\beta_{0j}.
\]
When $x=1$,
\[
\log\left(\frac{\pi_{j1}}{\pi_{11}}\right)
=
\beta_{0j}+\beta_{1j}.
\]
Here $\pi_{jk}=P(Y=j\mid x=k)$ for $k\in\{0,1\}$.

\paragraph{Odds Ratio}
The odds ratio comparing $x=1$ to $x=0$ for category $j$ relative to the
baseline category $1$ is
\[
OR_j
=
\frac{\pi_{j1}/\pi_{11}}{\pi_{j0}/\pi_{10}}.
\]
Using the model,
\[
\log(OR_j)
=
\log\left(\frac{\pi_{j1}}{\pi_{11}}\right)
-
\log\left(\frac{\pi_{j0}}{\pi_{10}}\right)
=
(\beta_{0j}+\beta_{1j})-\beta_{0j}
=
\beta_{1j}.
\]

Therefore,
\[
OR_j = e^{\beta_{1j}}.
\]

\paragraph{Estimated Odds Ratio}
With estimated probabilities $\hat\pi_{jk}$,
\[
\widehat{OR}_j
=
\frac{\hat\pi_{j1}/\hat\pi_{11}}
     {\hat\pi_{j0}/\hat\pi_{10}}
=
\exp(\hat\beta_{1j}).
\]
The CI is 
\[
\exp\!\left(
\hat{\beta}_{1j} \pm 1.96\,\sqrt{\widehat{\operatorname{Var}}(\hat{\beta}_{1j})}
\right)
\]
which is equivalent to
$$
\left(e^{\hat{\beta}_{1 j}-1.96 S E\left(\hat{\beta}_{1 j}\right)}, e^{\hat{\beta}_{1 j}+1.96 S E\left(\hat{\beta}_{1 j}\right)}\right)
$$
\subsection{Nominal and Ordinal Regression}
Let
\[
Y_i \in \{1,2,\dots,K\}, \qquad
\Pr(Y_i=j \mid x_i)=\pi_{ij}, \qquad
\sum_{j=1}^K \pi_{ij}=1.
\]
For nominal logistic regression, choose category \(K\) as the baseline. Then, for
\(j=1,\dots,K-1\),
\[
\log\left(\frac{\pi_{ij}}{\pi_{iK}}\right)=x_i^\top \beta_j.
\]
Equivalently,
\[
\pi_{ij}
=
\frac{\exp(x_i^\top \beta_j)}
{1+\sum_{l=1}^{K-1}\exp(x_i^\top \beta_l)},
\qquad j=1,\dots,K-1,
\]
and
\[
\pi_{iK}
=
\frac{1}
{1+\sum_{l=1}^{K-1}\exp(x_i^\top \beta_l)}.
\]
\[
Y_i \in \{1,2,\dots,K\}, \qquad
\Pr(Y_i \le j \mid x_i)=\gamma_{ij}, \qquad j=1,\dots,K-1.
\]
The ordinal logistic regression model (proportional odds model) is
\[
\log\left(\frac{\Pr(Y_i \le j \mid x_i)}{\Pr(Y_i > j \mid x_i)}\right)
=
\log\left(\frac{\gamma_{ij}}{1-\gamma_{ij}}\right)
=
\alpha_j - x_i^\top \beta,
\qquad j=1,\dots,K-1.
\]
Equivalently,
\[
\Pr(Y_i \le j \mid x_i)
=
\gamma_{ij}
=
\frac{\exp(\alpha_j - x_i^\top \beta)}
{1+\exp(\alpha_j - x_i^\top \beta)},
\qquad j=1,\dots,K-1.
\]

The category probabilities are
\[
\Pr(Y_i=1\mid x_i)=\gamma_{i1},
\]
\[
\Pr(Y_i=j\mid x_i)=\gamma_{ij}-\gamma_{i,j-1},
\qquad j=2,\dots,K-1,
\]
\[
\Pr(Y_i=K\mid x_i)=1-\gamma_{i,K-1}.
\]
\subsection*{Properties of the proportional odds model.}
For an ordinal response \(Y_i \in \{1,\dots,K\}\), define
\[
\gamma_{ij}=\Pr(Y_i \le j \mid x_i), \qquad j=1,\dots,K-1.
\]
The proportional odds model is
\[
\log\left(\frac{\gamma_{ij}}{1-\gamma_{ij}}\right)=\alpha_j-x_i^\top\beta.
\]

Its main properties are:
\begin{enumerate}
    \item It is a model for \emph{ordered} categorical responses.
    \item It models \emph{cumulative logits}.
    \item The slope vector \(\beta\) is the same for all cutpoints \(j\) (proportional odds assumption).
    \item Only the intercepts \(\alpha_j\) vary with \(j\).
    \item The cumulative logit curves are parallel in \(x_i\).
    \item The intercepts satisfy
    \[
    \alpha_1<\alpha_2<\cdots<\alpha_{K-1}.
    \]
    \item For a one-unit increase in \(x_r\), the cumulative odds are multiplied by
    \[
    \exp(-\beta_r),
    \]
    uniformly over all cutpoints.
    \item Equivalently, the odds of being in a higher category are multiplied by
    \[
    \exp(\beta_r).
    \]
    \item The category probabilities are recovered from the cumulative probabilities by
    \[
    \Pr(Y_i=1\mid x_i)=\gamma_{i1},
    \]
    \[
    \Pr(Y_i=j\mid x_i)=\gamma_{ij}-\gamma_{i,j-1}, \qquad j=2,\dots,K-1,
    \]
    \[
    \Pr(Y_i=K\mid x_i)=1-\gamma_{i,K-1}.
    \]
    \item The model uses fewer parameters than nominal logistic regression.
    \item When \(K=2\), it reduces to ordinary binary logistic regression.
    \item It is appropriate only when the proportional odds assumption is reasonable.
\end{enumerate}

\paragraph{Adjacent-categories logit model}
\[
\log\left(\frac{\pi_{ij}}{\pi_{i,j+1}}\right)=\alpha_j-x_i^\top\beta,
\qquad j=1,\dots,K-1.
\]
\textbf{Main properties:}
\begin{enumerate}
    \item It is for ordinal categorical responses.
    \item It models \emph{adjacent-category} odds, not cumulative odds.
    \item The same slope vector \(\beta\) is used for all adjacent logits.
    \item Only the intercepts \(\alpha_j\) vary with \(j\).
    \item For a one-unit increase in \(x_r\), the odds of category \(j\) versus \(j+1\) are multiplied by
    \[
    \exp(-\beta_r).
    \]
    \item It has \((K-1)+p\) parameters.
    \item It uses the ordering of the response categories.
    \item When \(K=2\), it reduces to binary logistic regression.
\end{enumerate}
\paragraph{Continuation-ratio logit model}
\[
\pi_{ij}=\Pr(Y_i=j\mid x_i), \qquad j=1,\dots,K.
\]
The continuation-ratio logit model is
\[
\log\left(\frac{\Pr(Y_i=j\mid x_i)}{\Pr(Y_i>j\mid x_i)}\right)
=
\log\left(\frac{\pi_{ij}}{\pi_{i,j+1}+\cdots+\pi_{iK}}\right)
=
\alpha_j-x_i^\top\beta,
\qquad j=1,\dots,K-1.
\]
Equivalently,
\[
\log\left(
\frac{\Pr(Y_i=j\mid Y_i\ge j,\;x_i)}
{\Pr(Y_i>j\mid Y_i\ge j,\;x_i)}
\right)
=
\alpha_j-x_i^\top\beta.
\]

\textbf{Main properties:}
\begin{enumerate}
    \item It is for ordinal categorical responses.
    \item It models the odds of category \(j\) versus all higher categories.
    \item It has a sequential interpretation.
    \item The same slope vector \(\beta\) is used for all logits.
    \item Only the intercepts \(\alpha_j\) vary with \(j\).
    \item For a one-unit increase in \(x_r\), the odds are multiplied by
    \[
    \exp(-\beta_r).
    \]
    \item It has \((K-1)+p\) parameters.
    \item It reduces to binary logistic regression when \(K=2\).
\end{enumerate}
Model can be compared with goodness of fit test
\newpage
\section{Quasi-likelihood for Generalized Linear Models}

Let $Y_1,\dots,Y_n$ be independent responses with
\[
\mathbb{E}(Y_i)=\mu_i,\qquad \mathrm{Var}(Y_i)=\phi V(\mu_i),
\]
 where $V(\mu_i)$ is the variance function and $\phi>0$ is the dispersion parameter.

In the quasi-likelihood approach, one does not assume a full probability model for $Y_i$. Instead, one seeks a function playing the role of a log-likelihood whose score with respect to $\mu_i$ has the same first two moments as the usual likelihood score in a correctly specified parametric model.

\subsection{Quasi-score conditions}
For each observation $i$, let $\ell_i(\mu_i;y_i)$ denote the contribution of the $i$th observation to a log-likelihood-type function. We require that its derivative with respect to $\mu_i$ satisfy the following three conditions:
\begin{align}
\mathbb{E}\!\left(\frac{\partial \ell_i}{\partial \mu_i}\right) &= 0, \label{eq:ql-cond1}\\
\mathrm{Var}\!\left(\frac{\partial \ell_i}{\partial \mu_i}\right) &= \frac{1}{\phi V(\mu_i)}, \label{eq:ql-cond2}\\
-\mathbb{E}\!\left(\frac{\partial^2 \ell_i}{\partial \mu_i^2}\right) &= \frac{1}{\phi V(\mu_i)}. \label{eq:ql-cond3}
\end{align}
Thus the score is unbiased, and its variance equals the expected information.

\begin{theorem}[Quasi-score identities]
Suppose
\[
\mathbb{E}(Y_i)=\mu_i,\qquad \mathrm{Var}(Y_i)=\phi V(\mu_i).
\]
Define
\[
q_i(\mu_i;y_i)=\frac{y_i-\mu_i}{\phi V(\mu_i)}.
\]
Then:
\begin{align*}
\mathbb{E}\{q_i(\mu_i;Y_i)\} &= 0,\\
\mathrm{Var}\{q_i(\mu_i;Y_i)\} &= \frac{1}{\phi V(\mu_i)},\\
-\mathbb{E}\!\left(\frac{\partial q_i(\mu_i;Y_i)}{\partial \mu_i}\right) &= \frac{1}{\phi V(\mu_i)}.
\end{align*}
Hence $q_i$ satisfies the quasi-likelihood score conditions.
\end{theorem}

\begin{proof}
First,
\[
\mathbb{E}\{q_i(\mu_i;Y_i)\}
=
\mathbb{E}\left(\frac{Y_i-\mu_i}{\phi V(\mu_i)}\right)
=
\frac{\mathbb{E}(Y_i)-\mu_i}{\phi V(\mu_i)}
=
0.
\]
Second,
\[
\mathrm{Var}\{q_i(\mu_i;Y_i)\}
=
\mathrm{Var}\left(\frac{Y_i-\mu_i}{\phi V(\mu_i)}\right)
=
\frac{\mathrm{Var}(Y_i)}{\phi^2 V(\mu_i)^2}
=
\frac{\phi V(\mu_i)}{\phi^2 V(\mu_i)^2}
=
\frac{1}{\phi V(\mu_i)}.
\]

Finally, differentiating $q_i(\mu_i;y_i)$ with respect to $\mu_i$ gives
\[
\frac{\partial q_i}{\partial \mu_i}
=
-\frac{1}{\phi V(\mu_i)}
-\frac{(y_i-\mu_i)V'(\mu_i)}{\phi V(\mu_i)^2}.
\]
Taking expectation,
\[
\mathbb{E}\!\left(\frac{\partial q_i}{\partial \mu_i}\right)
=
-\frac{1}{\phi V(\mu_i)}
-\frac{V'(\mu_i)}{\phi V(\mu_i)^2}\,
\mathbb{E}(Y_i-\mu_i).
\]
Since $\mathbb{E}(Y_i-\mu_i)=0$, this becomes
\[
\mathbb{E}\!\left(\frac{\partial q_i}{\partial \mu_i}\right)
=
-\frac{1}{\phi V(\mu_i)}.
\]
Therefore,
\[
-\mathbb{E}\!\left(\frac{\partial q_i}{\partial \mu_i}\right)
=
\frac{1}{\phi V(\mu_i)}.
\]
This proves the result.
\end{proof}
Since the quasi-score is defined by
\[
\frac{\partial Q_i}{\partial \mu_i}
=
\frac{y_i-\mu_i}{\phi V(\mu_i)},
\]
the associated quasi-likelihood contribution from the $i$th observation is obtained by integration:
\[
Q_i(\mu_i;y_i)
=
\int_{y_i}^{\mu_i}\frac{y_i-t}{\phi V(t)}\,dt,
\]
where $t$ is a dummy variable of integration.

Equivalently,
\[
Q_i(\mu_i;y_i)
=
-\int_{\mu_i}^{y_i}\frac{y_i-t}{\phi V(t)}\,dt.
\]
By construction,
\[
\frac{\partial Q_i(\mu_i;y_i)}{\partial \mu_i}
=
\frac{y_i-\mu_i}{\phi V(\mu_i)}.
\]
The overall quasi-likelihood is then
\[
Q(\mu;y)=\sum_{i=1}^n Q_i(\mu_i;y_i).
\]

\subsection*{Example: normal case}

Take
\[
\phi=\sigma^2=1,
\qquad
b(\theta(\mu_i))=\frac{\theta(\mu_i)^2}{2}.
\]
Then for the normal exponential family,
\[
\mu_i=b'(\theta_i)=\theta_i,
\qquad
V(\mu_i)=b''(\theta_i)=1.
\]
Hence
\[
\frac{\partial Q_i}{\partial \mu_i}
=
\frac{y_i-\mu_i}{\phi V(\mu_i)}
=
y_i-\mu_i.
\]
Integrating,
\[
Q_i(\mu_i;y_i)
=
\int_{y_i}^{\mu_i}(y_i-t)\,dt
=
-\frac{(y_i-\mu_i)^2}{2}.
\]

Thus the total quasi-likelihood is
\[
Q(\mu)
=
\sum_{i=1}^n Q_i
=
-\frac{1}{2}\sum_{i=1}^n (y_i-\mu_i)^2.
\]

If $\mu_i=\theta$ for all $i$, then
\[
Q(\theta)
=
-\frac{1}{2}\sum_{i=1}^n (y_i-\theta)^2.
\]
Maximizing $Q(\theta)$ is equivalent to minimizing
\[
\sum_{i=1}^n (y_i-\theta)^2.
\]
So
\[
\hat\theta=\bar y.
\]
It is consistent but less efficient
\[
\operatorname{Var}(\theta)\geq I(\theta)^{-1}
\]
\subsection{Quasi-deviance Computation}

The quasi-deviance is defined by
\[
D(y,\mu)
=
2\sum_{i=1}^n \{Q_i(y_i;y_i)-Q_i(\mu_i;y_i)\}.
\]
Equivalently, using the integral form of the quasi-likelihood,
\[
D(y,\mu)
=
2\sum_{i=1}^n \int_{\mu_i}^{y_i}\frac{y_i-t}{\phi V(t)}\,dt.
\]
\textbf{In the normal case} with
\[
\phi=1,\qquad V(\mu_i)=1,
\]
we have
\[
Q_i(\mu_i;y_i)=-\frac{(y_i-\mu_i)^2}{2},
\qquad
Q_i(y_i;y_i)=0.
\]
Hence
\[
D(y,\mu)
=
2\sum_{i=1}^n \left(0+\frac{(y_i-\mu_i)^2}{2}\right)
=
\sum_{i=1}^n (y_i-\mu_i)^2.
\]
So for the Gaussian model, the quasi-deviance is just the residual sum of squares.

\subsection{$\beta$ estimation in a GLM via quasi-likelihood}

In a GLM, let
\[
g(\mu_i)=\eta_i=x_i^T\beta,
\]
where $g(\cdot)$ is the link function. The quasi-likelihood estimator $\hat\beta$ is obtained by solving
\[
\frac{\partial Q}{\partial \beta}
=
\sum_{i=1}^n
\frac{\partial Q(\mu_i,y_i)}{\partial \mu_i}
\frac{\partial \mu_i}{\partial \beta}
=
0.
\]
Since
\[
\frac{\partial Q(\mu_i,y_i)}{\partial \mu_i}
=
\frac{y_i-\mu_i}{\phi V(\mu_i)},
\]
we get the estimating equation
\[
\sum_{i=1}^n
\frac{y_i-\mu_i}{\phi V(\mu_i)}
\frac{\partial \mu_i}{\partial \beta}
=
0.
\]

Now
\[
\frac{\partial \mu_i}{\partial \beta}
=
\frac{d\mu_i}{d\eta_i}\frac{\partial \eta_i}{\partial \beta}
=
\frac{d\mu_i}{d\eta_i}x_i,
\]
so the quasi-score for $\beta$ is
\[
U(\beta)
=
\sum_{i=1}^n
\frac{(y_i-\mu_i)}{\phi V(\mu_i)}
\frac{1}{g^\prime(\mu_i)}x_i
=
0.
\]
In matrix form, this is
\[
U(\beta)
=
\frac{D^TV^{-1}(y-\mu)}{\sigma^2}=0,
\]
where
\[
V = diag(V(\mu_1),\dots,V(\mu_n)),\quad D = diag(\frac{\partial \mu_1}{\partial \beta},\dots,(\frac{\partial \mu_n}{\partial \beta})
\]
Thus $\hat\beta$ is the solution of the quasi-score equation
\[
\sum_{i=1}^n
\frac{(y_i-\mu_i)}{V(\mu_i)}
\frac{d\mu_i}{d\eta_i}x_i
=
0,
\]
up to the constant factor $1/\phi$.
\subsection{Newton--Raphson to solve Quasi-Score GLM}

Write the quasi-score as
\[
U(\beta)
=
\frac{D^T V^{-1}(y-\mu)}{\sigma^2}
=
0,
\]
where
\[
V=\operatorname{diag}\bigl(V(\mu_1),\dots,V(\mu_n)\bigr),
\]
and
\[
D=\frac{\partial \mu}{\partial \beta^T}
=
\begin{pmatrix}
\frac{\partial \mu_1}{\partial \beta_1} & \cdots & \frac{\partial \mu_1}{\partial \beta_p}\\
\vdots & \ddots & \vdots\\
\frac{\partial \mu_n}{\partial \beta_1} & \cdots & \frac{\partial \mu_n}{\partial \beta_p}
\end{pmatrix}.
\]
Thus $D$ is the Jacobian matrix of $\mu=(\mu_1,\dots,\mu_n)^T$ with respect to $\beta$.

Newton--Raphson updates $\beta$ by
\[
\beta^{(m+1)}
=
\beta^{(m)}
-
\left[
\frac{\partial U(\beta)}{\partial \beta^T}
\right]^{-1}_{\beta=\beta^{(m)}}
U(\beta^{(m)}).
\]
Using the usual Fisher scoring approximation,
\[
-\frac{\partial U(\beta)}{\partial \beta^T}
\approx
\frac{D^T V^{-1} D}{\sigma^2},
\]
so the update becomes
\[
\beta^{(m+1)}
=
\beta^{(m)}
+
(D^T V^{-1}D)^{-1}D^T V^{-1}(y-\mu).
\]
Equivalently,
\[
\Delta\beta
=
(D^T V^{-1}D)^{-1}D^T V^{-1}(y-\mu),
\qquad
\beta^{(m+1)}=\beta^{(m)}+\Delta\beta.
\]
This is the Newton--Raphson/Fisher scoring step for estimating $\beta$ in the GLM quasi-likelihood framework.
\[
U(\beta)=\frac{D^T V^{-1}(y-\mu)}{\sigma^2}.
\]
Since
\[
\mathrm{Cov}(y)=\sigma^2 V,
\]
we have
\[
\mathrm{Cov}\{U(\beta)\}
=
\frac{1}{\sigma^4}D^T V^{-1}\,\mathrm{Cov}(y)\,V^{-1}D
=
\frac{1}{\sigma^4}D^T V^{-1}(\sigma^2 V)V^{-1}D.
\]
Therefore,
\[
\mathrm{Cov}\{U(\beta)\}
=
\frac{1}{\sigma^2}D^T V^{-1}D.
\]
\newpage
\section{Generalized Estimating Equations}
For clustered/longitudinal data, let there be \(n\) independent subjects (clusters), where subject \(i\) has \(m_i\) repeated measurements.

The response vector for subject \(i\) is
\[
Y_i =
\begin{pmatrix}
Y_{i1}\\
Y_{i2}\\
\vdots\\
Y_{im_i}
\end{pmatrix},
\qquad i=1,\dots,n.
\]
The corresponding covariate matrix is
\[
X_i =
\begin{pmatrix}
X_{i1}^\top\\
X_{i2}^\top\\
\vdots\\
X_{im_i}^\top
\end{pmatrix},
\]
where each \(X_{ij}\) is a \(p\times 1\) covariate vector, so that
\[
X_i \in \mathbb{R}^{m_i\times p},
\qquad
Y_i \in \mathbb{R}^{m_i}.
\]
The marginal mean vector is
\[
\mu_i = \mathbb{E}(Y_i \mid X_i)
=
\begin{pmatrix}
\mu_{i1}\\
\mu_{i2}\\
\vdots\\
\mu_{im_i}
\end{pmatrix}.
\]
The GEE marginal model specifies
\[
g(\mu_{ij}) =\eta_{ij}= X_{ij}^\top \beta,
\qquad j=1,\dots,m_i,
\]
or equivalently, in vector form,
\[
g(\mu_i)=
\begin{pmatrix}
g(\mu_{i1})\\
g(\mu_{i2})\\
\vdots\\
g(\mu_{im_i})
\end{pmatrix}
= X_i \beta.
\]
The marginal variance is written as
\[
\mathrm{Var}(Y_i\mid X_i)=V_i
=
A_i^{1/2} R_i(\alpha) A_i^{1/2},
\]
where
\[
A_i=\mathrm{diag}\!\bigl(\phi V(\mu_{i1}),\dots,\phi V(\mu_{im_i})\bigr)
\]
contains the marginal variances, and \(R_i(\alpha)\) is the working correlation matrix for the repeated measurements within subject \(i\).

For the marginal model
\[
g(\mu_{ij}) = X_{ij}^\top \beta,
\qquad
\mu_i = \mathbb{E}(Y_i\mid X_i),
\]
the generalized estimating equation (GEE) for \(\beta\) is
\[
U(\beta)
=
\sum_{i=1}^n D_i^\top V_i^{-1}(Y_i-\mu_i)=0,
\]
where
\[
D_i=\frac{\partial \mu_i}{\partial \beta^\top},
\qquad
V_i=\mathrm{Var}(Y_i\mid X_i).
\]
Using the GEE working covariance structure,
\[
V_i=A_i^{1/2}R_i(\alpha)A_i^{1/2},
\]
with
\[
A_i=\mathrm{diag}\!\bigl(\phi V(\mu_{i1}),\dots,\phi V(\mu_{im_i})\bigr),
\]
we get
\[
\sum_{i=1}^n
D_i^\top
A_i^{-1/2}R_i(\alpha)^{-1}A_i^{-1/2}
(Y_i-\mu_i)=0.
\]
This is the extension of the GLM quasi-score equation
\[
\sum_{i=1}^n D_i^\top V_i^{-1}(Y_i-\mu_i)=0
\]
from independent data to clustered/longitudinal data by incorporating the within-cluster working correlation \(R_i(\alpha)\).

\subsection{Solving GEE}
A common two-step iterative algorithm for solving the GEE is:
\[
U(\beta)
=
\sum_{i=1}^n D_i^\top V_i^{-1}(Y_i-\mu_i)=0,
\qquad
V_i=A_i^{1/2}R_i(\alpha)A_i^{1/2}.
\]
\textbf{Step 1: Update \(\beta\) given \(\alpha,\phi\).}  

Given current estimates \(\alpha^{(t)},\phi^{(t)}\), form
\[
V_i^{(t)}=A_i^{(t)1/2}R_i\!\left(\alpha^{(t)}\right)A_i^{(t)1/2},
\]
and solve
\[
\sum_{i=1}^n D_i^\top \bigl(V_i^{(t)}\bigr)^{-1}(Y_i-\mu_i)=0
\]
for \(\beta\), usually by Fisher scoring / Newton--Raphson:
\[
\beta^{(t+1)}
=
\beta^{(t)}
+
\left[
\sum_{i=1}^n D_i^\top \bigl(V_i^{(t)}\bigr)^{-1} D_i
\right]^{-1}
\sum_{i=1}^n D_i^\top \bigl(V_i^{(t)}\bigr)^{-1}(Y_i-\mu_i).
\]
\textbf{Step 2: Update \(\alpha\) and \(\phi\) given \(\beta\).}  

Using the updated \(\beta^{(t+1)}\), compute fitted means
\[
\mu_{ij}^{(t+1)}=g^{-1}(X_{ij}^\top \beta^{(t+1)}),
\]
then update the nuisance parameters:
\begin{itemize}
\item estimate \(\phi^{(t+1)}\) from Pearson residuals;
\item estimate \(\alpha^{(t+1)}\) from the within-cluster residual correlations.
\end{itemize}
Repeat Steps 1 and 2 until convergence of \(\beta\) (and usually also \(\alpha,\phi\)).
\subsection*{GEE estimation consistency}
Under regularity conditions, the GEE estimator \(\hat\beta\) is consistent and asymptotically normal:
\[
\hat\beta \xrightarrow{p} \beta_0,
\]
and
\[
\sqrt{n}\,(\hat\beta-\beta_0)
\xrightarrow{d}
N\!\left(0,\; A^{-1} B A^{-1}\right),
\]
where
\[
A
=
\lim_{n\to\infty}
\frac{1}{n}
\sum_{i=1}^n
D_i^\top V_i^{-1} D_i,
\]
and
\[
B
=
\lim_{n\to\infty}
\frac{1}{n}
\sum_{i=1}^n
D_i^\top V_i^{-1}
\mathrm{Cov}(Y_i)
V_i^{-1} D_i.
\]
Thus the asymptotic covariance matrix is the \emph{sandwich form}
\[
\mathrm{Var}(\hat\beta)
\approx
\frac{1}{n}A^{-1}BA^{-1}.
\]
If the working correlation is correctly specified, then \(B=A\), so
\[
\sqrt{n}\,(\hat\beta-\beta_0)
\xrightarrow{d}
N\!\left(0,\; A^{-1}\right).
\]
\subsection{Covariance Matrix Estimation}
A consistent estimator of the asymptotic covariance is the robust sandwich estimator. That can be estimated with:
\[
\widehat{\mathrm{Var}}(\hat\beta)
=
\left(\sum_{i=1}^n D_i^\top V_i^{-1} D_i\right)^{-1}
\left(\sum_{i=1}^n D_i^\top V_i^{-1}(Y_i-\mu_i)(Y_i-\mu_i)^\top V_i^{-1} D_i\right)
\left(\sum_{i=1}^n D_i^\top V_i^{-1} D_i\right)^{-1}.
\]
For the \emph{compound symmetry} (exchangeable) working correlation structure, the within-cluster correlation matrix is
\[
R_i(\alpha)=
\begin{pmatrix}
1 & \alpha & \alpha & \cdots & \alpha\\
\alpha & 1 & \alpha & \cdots & \alpha\\
\alpha & \alpha & 1 & \cdots & \alpha\\
\vdots & \vdots & \vdots & \ddots & \vdots\\
\alpha & \alpha & \alpha & \cdots & 1
\end{pmatrix}_{m_i\times m_i}.
\]
Thus every pair of distinct observations in the same cluster has the same correlation:
\[
\mathrm{Corr}(Y_{ij},Y_{ik})=\alpha,\qquad j\neq k.
\]
The corresponding working covariance matrix is
\[
V_i
=
A_i^{1/2}R_i(\alpha)A_i^{1/2}.
\]
If all marginal variances are equal to \(\sigma^2\), then
\[
V_i=\sigma^2 R_i(\alpha).
\]
Under the compound symmetry (exchangeable) working correlation structure, define the Pearson residual by
\[
r_{ij}
=
\frac{Y_{ij}-\hat\mu_{ij}}
{\sqrt{\hat\phi\,V(\hat\mu_{ij})}}.
\]
Then the common correlation parameter \(\alpha\) is typically estimated by the average of the within-cluster Pearson residual cross-products:
\[
\hat\alpha
=
\frac{\displaystyle \sum_{i=1}^n \sum_{j<k} r_{ij}r_{ik}}
{\displaystyle \sum_{i=1}^n \binom{m_i}{2}}.
\]
Equivalently, writing it in terms of \(Y_{ij}\) and \(\hat\mu_{ij}\),
\[
\hat\alpha
=
\frac{\displaystyle \sum_{i=1}^n \sum_{j<k}
\frac{(Y_{ij}-\hat\mu_{ij})(Y_{ik}-\hat\mu_{ik})}
{\sqrt{\hat\phi\,V(\hat\mu_{ij})}\sqrt{\hat\phi\,V(\hat\mu_{ik})}}}
{\displaystyle \sum_{i=1}^n \binom{m_i}{2}}.
\]
So for exchangeable correlation, \(\hat\alpha\) is just the average empirical Pearson correlation over all pairs within the same cluster.

\subsection*{Example: Gaussian Case}
In the Gaussian case, suppose
\[
Y_i \sim N(\mu_i,V_i), \qquad \mu_i = X_i\beta,
\]
with clusters independent, i.e.
\[
\mathrm{Cov}(Y_i,Y_{i'})=0, \qquad i\neq i'.
\]
If each \(V_i\) is known, then the GEE becomes
\[
\sum_{i=1}^n D_i^\top V_i^{-1}(Y_i-\mu_i)=0.
\]
Since
\[
\mu_i=X_i\beta,
\qquad
D_i=\frac{\partial \mu_i}{\partial \beta^\top}=X_i,
\]
the estimating equation is
\[
\sum_{i=1}^n X_i^\top V_i^{-1}(Y_i-X_i\beta)=0.
\]
Solving for \(\beta\) gives the generalized least squares estimator
\[
\hat\beta
=
\left(\sum_{i=1}^n X_i^\top V_i^{-1}X_i\right)^{-1}
\left(\sum_{i=1}^n X_i^\top V_i^{-1}Y_i\right).
\]
Equivalently, stacking all clusters,
\[
Y=
\begin{pmatrix}
Y_1\\ \vdots\\ Y_n
\end{pmatrix},
\qquad
X=
\begin{pmatrix}
X_1\\ \vdots\\ X_n
\end{pmatrix},
\qquad
V=\mathrm{blockdiag}(V_1,\dots,V_n),
\]
we have
\[
\hat\beta=(X^\top V^{-1}X)^{-1}X^\top V^{-1}Y.
\]
Its covariance matrix is
\[
\mathrm{Var}(\hat\beta)=(X^\top V^{-1}X)^{-1}.
\]
If the covariance matrices \(V_i\) are unknown, a common approach is \emph{iterative feasible generalized least squares}.
\medskip

\textbf{Step 1: Initial estimate of \(\beta\).}  
Start with an initial working covariance, for example
\[
V_i^{(0)}=I_{m_i},
\]
and compute the OLS estimator
\[
\beta^{(0)}
=
\left(\sum_{i=1}^n X_i^\top X_i\right)^{-1}
\left(\sum_{i=1}^n X_i^\top Y_i\right).
\]
\medskip
\textbf{Step 2: Estimate \(V_i\) from residuals.}  
Using the current estimate \(\beta^{(t)}\), form residuals
\[
e_i^{(t)} = Y_i - X_i\beta^{(t)}.
\]
Then estimate the covariance matrix \(V_i\) by a parametric working model
\[
V_i = V_i(\theta),
\]
with parameter \(\theta\) estimated from the residuals \(e_i^{(t)}\). For example, estimate variance and correlation parameters by moment methods or maximum likelihood.
\medskip

\textbf{Step 3: Update \(\beta\).}  
Given the updated covariance estimate \(\hat V_i^{(t)}\), solve
\[
\sum_{i=1}^n X_i^\top \bigl(\hat V_i^{(t)}\bigr)^{-1}(Y_i-X_i\beta)=0,
\]
so that
\[
\beta^{(t+1)}
=
\left(\sum_{i=1}^n X_i^\top \bigl(\hat V_i^{(t)}\bigr)^{-1}X_i\right)^{-1}
\left(\sum_{i=1}^n X_i^\top \bigl(\hat V_i^{(t)}\bigr)^{-1}Y_i\right).
\]

\medskip
\textbf{Step 4: Iterate.}  
Repeat Steps 2 and 3 until convergence of \(\beta\) and the covariance parameters.
\medskip
Thus the iterative scheme is
\[
\beta^{(t)}
\longrightarrow
e_i^{(t)}=Y_i-X_i\beta^{(t)}
\longrightarrow
\hat V_i^{(t)}
\longrightarrow
\beta^{(t+1)}.
\]
This is the usual \emph{feasible GLS} / \emph{iterative GEE} procedure in the Gaussian case.
\subsection*{The choice of $V$ estimation}
If \(\hat V\) is substituted for \(V\) in the estimation of
\[
\mathrm{Var}(\hat\beta)=(X^\top V^{-1}X)^{-1},
\]
then the resulting estimator
\[
(X^\top \hat V^{-1}X)^{-1}
\]
usually underestimates the true variability of \(\hat\beta\), because it treats the estimated covariance matrix \(\hat V\) as if it were known and therefore ignores the additional variability from estimating \(V\).

When \(V_i\) is replaced by \(\hat V_i\), this plug-in variance formula does not account for the uncertainty in estimating the covariance parameters. As a result, it is typically too small.

A more reliable large-sample variance estimator is the sandwich estimator:
\[
\widehat{\mathrm{Var}}(\hat\beta)
=
A^{-1}BA^{-1},
\]
where
\[
A=\sum_{i=1}^n X_i^\top \hat V_i^{-1}X_i,
\]
and
\[
B=\sum_{i=1}^n X_i^\top \hat V_i^{-1}
(Y_i-X_i\hat\beta)(Y_i-X_i\hat\beta)^\top
\hat V_i^{-1}X_i.
\]

Thus,
\[
\widehat{\mathrm{Var}}(\hat\beta)
=
\left(\sum_{i=1}^n X_i^\top \hat V_i^{-1}X_i\right)^{-1}
\left(\sum_{i=1}^n X_i^\top \hat V_i^{-1}
(Y_i-X_i\hat\beta)(Y_i-X_i\hat\beta)^\top
\hat V_i^{-1}X_i\right)
\left(\sum_{i=1}^n X_i^\top \hat V_i^{-1}X_i\right)^{-1}.
\]

This estimator remains consistent even if the working covariance structure is misspecified, and it corrects for the downward bias of the naive plug-in estimator. It is consistent and robust to the misspecification of $V$.

\subsection*{Small-Sample Variance Estimation in GEE}
Although the sandwich estimator is consistent as the number of independent
clusters \(K\) tends to infinity, it tends to underestimate the true
variance when the number of clusters is small. In particular, when
\(K\leq 30\), the resulting standard errors may be too small and Wald tests
may have inflated Type I error rates.

A cluster-level jackknife variance estimator can be used to reduce this
small-sample bias. Let
\[
\widehat{\boldsymbol{\beta}}_{-i}
\]
denote the GEE estimator obtained after deleting the entire \(i\)-th cluster.
The delete-one-cluster jackknife variance estimator is
\[
\widehat{\operatorname{Var}}_{\mathrm{JK}}
\left(
\widehat{\boldsymbol{\beta}}
\right)
=
\frac{K-1}{K}
\sum_{i=1}^{K}
\left(
\widehat{\boldsymbol{\beta}}_{-i}
-
\overline{\boldsymbol{\beta}}_{(-)}
\right)
\left(
\widehat{\boldsymbol{\beta}}_{-i}
-
\overline{\boldsymbol{\beta}}_{(-)}
\right)^{T},
\]
where
\[
\overline{\boldsymbol{\beta}}_{(-)}
=
\frac{1}{K}
\sum_{i=1}^{K}
\widehat{\boldsymbol{\beta}}_{-i}.
\]

An alternative finite-sample version, sometimes used in GEE applications,
centers the deleted-cluster estimates at the full-sample estimate and applies
a degrees-of-freedom correction:
\[
\widehat{\operatorname{Var}}_{\mathrm{JK}}
\left(
\widehat{\boldsymbol{\beta}}
\right)
=
\frac{K-p}{K}
\sum_{i=1}^{K}
\left(
\widehat{\boldsymbol{\beta}}_{-i}
-
\widehat{\boldsymbol{\beta}}
\right)
\left(
\widehat{\boldsymbol{\beta}}_{-i}
-
\widehat{\boldsymbol{\beta}}
\right)^{T},
\]
where \(p\) is the number of estimated regression parameters. The exact
scaling convention may differ across references and software packages.

The jackknife standard error for the \(j\)-th regression coefficient is
\[
\operatorname{SE}_{\mathrm{JK}}
\left(
\widehat{\beta}_j
\right)
=
\sqrt{
\left[
\widehat{\operatorname{Var}}_{\mathrm{JK}}
\left(
\widehat{\boldsymbol{\beta}}
\right)
\right]_{jj}
}.
\]

For a small number of clusters, inference is commonly based on a
\(t\)-distribution rather than the standard normal distribution:
\[
T_j
=
\frac{
\widehat{\beta}_j-\beta_{j,0}
}{
\operatorname{SE}_{\mathrm{JK}}
\left(
\widehat{\beta}_j
\right)
}
\overset{\cdot}{\sim}
t_{K-p}.
\]
An approximate \(100(1-\alpha)\%\) confidence interval is therefore
\[
\widehat{\beta}_j
\pm
t_{K-p,\,1-\alpha/2}
\operatorname{SE}_{\mathrm{JK}}
\left(
\widehat{\beta}_j
\right).
\]
Thus, the ordinary sandwich estimator is generally appropriate when the
number of independent clusters is large, whereas a jackknife or another
small-sample corrected variance estimator is preferable when \(K\) is small.

\subsection*{R Implementation for GEE}

The \texttt{geepack} package can be used to fit generalized estimating
equations through \texttt{geeglm()}.

\begin{verbatim}
library(geepack)

fit <- geeglm(
  y ~ x1 + x2,
  id = cluster_id,
  data = dat,
  family = binomial(),
  corstr = "exchangeable",
  std.err = "fij"
)
\end{verbatim}

The main standard-error options are
\[
\begin{array}{ll}
\texttt{san.se} & \text{ordinary sandwich standard errors},\\
\texttt{j1s}    & \text{one-step jackknife standard errors},\\
\texttt{jack}   & \text{approximate jackknife standard errors},\\
\texttt{fij}    & \text{fully iterated jackknife standard errors}.
\end{array}
\]

For a small number of clusters, \texttt{std.err = "fij"} is generally
preferred.

Small-sample corrected inference can also be performed using the
\texttt{clubSandwich} package:

\begin{verbatim}
library(clubSandwich)

coef_test(
  fit,
  vcov = "CR2",
  cluster = dat$cluster_id,
  test = "Satterthwaite"
)
\end{verbatim}

Thus, a practical choice is
\[
\boxed{
\texttt{geepack::geeglm()}
\text{ for GEE fitting, and }
\texttt{clubSandwich::coef\_test()}
\text{ for small-sample corrected inference.}
}
\]
\end{document}